% EHB415 Gorev Kartlari 2 (ikinci faz). Derleme: tectonic -X compile gorev_kartlari_2.tex  (fontlar\ klasoru gerekir)
\documentclass[10.5pt,a4paper]{article}
\usepackage{fontspec}
\setmainfont{georgia}[Path=./fontlar/,Extension=.ttf,BoldFont=georgiab,ItalicFont=georgiai,BoldItalicFont=georgiaz]
\setsansfont{arial}[Path=./fontlar/,Extension=.ttf,BoldFont=arialbd,ItalicFont=ariali,BoldItalicFont=arialbi]
\setmonofont{consola}[Path=./fontlar/,Extension=.ttf,BoldFont=consolab,ItalicFont=consolai,BoldItalicFont=consolaz,Scale=0.88]
\usepackage[a4paper,left=2cm,right=2cm,top=1.9cm,bottom=1.9cm]{geometry}
\usepackage{xcolor,booktabs,enumitem,amsmath,url,float}
\usepackage[most]{tcolorbox}
\usepackage{xltabular,longtable}
\usepackage{fancyhdr,titlesec,listings}
\usepackage[colorlinks=true,linkcolor=blue!50!black,urlcolor=blue!50!black]{hyperref}
\emergencystretch=2em
\setlength{\parindent}{0pt}\setlength{\parskip}{4pt plus 1pt}
\setlist{itemsep=2pt,topsep=3pt,parsep=0pt}
\definecolor{ink}{HTML}{1F2937}\definecolor{acc}{HTML}{1D4ED8}\definecolor{accl}{HTML}{DBEAFE}
\definecolor{warn}{HTML}{EA580C}\definecolor{warnl}{HTML}{FFEDD5}\definecolor{good}{HTML}{16A34A}
\definecolor{goodl}{HTML}{DCFCE7}\definecolor{muted}{HTML}{6B7280}\definecolor{codebg}{HTML}{F3F4F6}
\color{ink}
\titleformat{\section}{\sffamily\bfseries\Large\color{acc}}{}{0pt}{}[{\color{acc!35}\titlerule[1pt]}]
\titleformat{\subsection}{\sffamily\bfseries\normalsize\color{acc!80!black}}{}{0pt}{}
\titlespacing*{\section}{0pt}{6pt}{6pt}\titlespacing*{\subsection}{0pt}{8pt}{2pt}
\pagestyle{fancy}\fancyhf{}\renewcommand{\headrulewidth}{0.3pt}
\fancyhead[L]{\sffamily\footnotesize\color{muted} EHB415 \textbullet\ ISAC projesi \textbullet\ Görev kartları 2}
\fancyfoot[C]{\sffamily\footnotesize\color{muted}\thepage}
\lstset{basicstyle=\ttfamily\footnotesize,backgroundcolor=\color{codebg},frame=single,rulecolor=\color{codebg!70!black},
  breaklines=true,columns=fullflexible,keepspaces=true,showstringspaces=false,aboveskip=4pt,belowskip=6pt,
  literate={ğ}{g}1 {ş}{s}1 {ı}{i}1 {ü}{u}1 {ö}{o}1 {ç}{c}1}
\newcommand{\kod}[1]{{\small\path{#1}}}
\newenvironment{komut}{\par\noindent\begin{minipage}{\textwidth}}{\end{minipage}\par}
\newtcolorbox{neden}{enhanced,colback=accl,colframe=acc,boxrule=0.5pt,arc=2pt,left=6pt,right=6pt,top=4pt,bottom=4pt,
  fonttitle=\sffamily\bfseries\small,title=Bu görev neden önemli?,coltitle=white,colbacktitle=acc,
  attach boxed title to top left={yshift=-2mm,xshift=4mm},boxed title style={arc=2pt,boxrule=0pt}}
\newtcolorbox{kontrol}{enhanced,colback=goodl,colframe=good,boxrule=0.5pt,arc=2pt,left=6pt,right=6pt,top=4pt,bottom=4pt,
  fonttitle=\sffamily\bfseries\small,title=Kontrol: bizde ölçülen sonuç (kendi sonucunla karşılaştır),coltitle=white,colbacktitle=good,
  attach boxed title to top left={yshift=-2mm,xshift=4mm},boxed title style={arc=2pt,boxrule=0pt}}
\newtcolorbox{teslim}{enhanced,colback=warnl,colframe=warn,boxrule=0.5pt,arc=2pt,left=6pt,right=6pt,top=4pt,bottom=4pt,
  fonttitle=\sffamily\bfseries\small,title=Ne teslim edeceksin?,coltitle=white,colbacktitle=warn,
  attach boxed title to top left={yshift=-2mm,xshift=4mm},boxed title style={arc=2pt,boxrule=0pt}}
\begin{document}

% ======================================================================= ORTAK
\section*{EHB415 ISAC projesi: görev kartları 2 (ikinci faz)}
{\sffamily\color{muted} Birinci fazda üç aracın kurulumunu doğruladık (birleşik rapor). Bu fazda her araçtan \emph{sayı} üretiyoruz. Önce Ortak bilgi'yi, sonra kendi kartını oku.}

\subsection*{Proje sorusu ve ortak eksen}
\textbf{Soru:} sensing'e ayrılan kaynak payı $\alpha$ arttıkça sensing ne kazanır, haberleşme ne kaybeder?
Sionna kazancı, ns-3 ve OCUDU kaybı ölçer. Ortak eksen $\alpha$ (sensing'e giden kaynağın toplam kaynağa oranı). Standartta (TR 38.765 §6.2) kaynak oranı Tip~1/2/3 ile tanımlıdır; Dickerson vd. modeli: $\alpha_{sense}=\frac{N_b}{1-OH}\cdot\frac{f_{sense}}{N_{sym/sn}}\cdot\frac{N_{PRB,sense}}{N_{PRB,toplam}}$ ve $T(\eta)=\min\{\eta,1-\alpha_{sense}\}$.

\subsection*{Kartlardaki işaretler ve kurallar}
\begin{itemize}[leftmargin=1.4em]
\item \textbf{[Okundu]}: o bölümü kartı hazırlayan kişi PDF'ten ya da kaynak koddan okudu; kartta yazanlar o metinde ya da kodda görülenlerdir.
\item \textbf{[Yalnızca özet]}: yalnızca özetini gördük; tam metni sen okuyacaksın, kartta onun hakkında özette yazandan fazlası yok.
\item \textbf{Kontrol kutusu}: komutu bu bilgisayarda çalıştırıp gördüğümüz sayılar. Sende farklı çıkarsa \textbf{farkı yaz}; hata sayılmaz.
\item Teslimde \textbf{çalıştırdığın komutu} ve \textbf{gördüğün çıktıyı} yaz. Anlamadığın yeri ``anlamadığım yer'' diye yaz.
\item Yeni dosyalar için: \texttt{git pull}. Yeni betikler: \kod{ogretici/betikler/ns3\_yuk\_sweep.sh}, \kod{ns3\_bant\_sweep.sh}, \kod{ocudu\_prs\_kos.sh} (yeni 5.\ argüman), \kod{ogretici/sionna\_sehir/tr\_metrik\_olc.py}; veriler \kod{ogretici/veri/}.
\end{itemize}

\begingroup\sffamily\small
\begin{xltabular}{\textwidth}{@{}p{1.5cm}p{2.6cm}Xp{3.6cm}@{}}
\toprule
\textbf{Kim} & \textbf{Alan} & \textbf{Bu fazda} & \textbf{Çıktı} \\
\midrule
\textbf{Taha} & 3GPP, metrik, rapor & TR 38.765 §6.3 sonuçları; $\alpha$ modelini OCUDU ve ns-3 ölçümleriyle karşılaştırma & 3 tablo + rapor bölümleri \\
\textbf{Nihat} & Sionna & TR metrikleri (Pf1, Pmd hedefi), $X_{SI}$ deneyi, varsayım karşılaştırması & Çıktı tablosu + karşılaştırma tablosu \\
\textbf{Aziz} & ns-3 / 5G-LENA & Yük ve bant taraması, Sionna-RT örneğinin durumu, ``hangi yol'' kararı & 2 tablo + 1 sayfa karar notu \\
\textbf{Fatih} & OCUDU & PRS ofset deneyi, kod okuma, çok kullanıcılı deney & Tablo + 2 sorunun cevabı \\
\bottomrule
\end{xltabular}
\endgroup
\textbf{Teslim tarihi:} \underline{\hspace{4cm}} \quad (kartı verirken yaz)

% ======================================================================= TAHA
\clearpage
\section*{Taha: 3GPP sonuçları, metrik ve $\alpha$ doğrulaması}
\begin{neden}
Rapor iki şeyi gösterecek: standardın kaynak oranı ve sonuçları nasıl verdiği, ve bizim $\alpha$ tanımımızın ölçümle tutup tutmadığı. TR 38.765 sonuçlarını \textbf{sensing kaynak oranına göre iki sınıfa} ($\le\%10$ ve $>\%10$) ayırıyor (Tablo 6.3-1 ve 6.3-2); bizim $\alpha$ ekseni tam bu sınıflandırmaya oturuyor.
\end{neden}

\subsection*{Okuma}
\begingroup\sffamily\small
\begin{xltabular}{\textwidth}{@{}p{3.0cm}p{3.3cm}Xp{1.6cm}@{}}
\toprule
\textbf{Kaynak} & \textbf{Yer} & \textbf{Ne arayacaksın} & \textbf{Durum} \\
\midrule
3GPP TR 38.765 v20.0.0 (\kod{makaleler/}) & §6.2 Tablo 6.2-1; §6.3.0; §6.3.2 Tablo 6.3-2 ve Case 2-1\ldots 2-6 (sayfa 15--20) & Baseline~2 (37\,dBm) bizim senaryomuz. Tablo 6.3-2: 6 durum; sütunlar: \emph{kaç TRP}, \emph{sensing kaynak oranı} ($\le\%10$ / $>\%10$), \emph{Tx ışın sayısı}. Her durumda $P_{md}$, $P_{f1}$, $P_{f2}$, konum ve hız aralıklarını not et. Tek TRP = bizim monostatik gNB. & [Okundu] \\
 & §5.1 & Ölçüm raporu seviyeleri A, B, C, D (ham veriden nesne düzeyine); her seviyenin neyi taşıdığı. & [Okundu] \\
\midrule
Dickerson vd. (arXiv 2606.21677) & III.B (denklem 9--12), IV.A, IV.D & IV.D: iki sabit politika: sensing-öncelikli \%15 ayırır (throughput sınırı 0,85), haberleşme-öncelikli \%3 ayırır (sınır 0,97). Bu iki sayı, $T=1-\alpha$ kuralının makaledeki örneği. & [Okundu] \\
\midrule
Blandino vd., \emph{Evaluation of gNB Monostatic Sensing for UAV Use Case} (2026, arXiv; Consensus ile bulundu) & Özet & 3GPP UMa-AV varsayımlarıyla gNB monostatic sensing; özet: PRS ve CP-OFDM kullanıyor, \%5'ten az yanlış alarmla \%70'ten fazla tespit, 90. yüzdelik hata dikeyde 4\,m, yatayda 6\,m; simülatör \texttt{5GNRad} yayınlanmış. TR 38.765 sonuçlarıyla aynı senaryo olduğu için raporda karşılaştırma kaynağı. & [Yalnızca özet] \\
\bottomrule
\end{xltabular}
\endgroup

\subsection*{Adımlar}
\begin{enumerate}[leftmargin=1.6em]
\item \textbf{TR Baseline~2 sonuç tablosu} (\kod{docs/metrics.md}'ye yeni bölüm). Satırlar Case 2-1\ldots 2-6; sütunlar: TRP sayısı, kaynak oranı, ışın, kaynak sayısı, $P_{md}$ aralığı, $P_{f2}$ aralığı, hedefler tuttu mu (evet / hayır). Değerleri \S6.3.2'den kendin yaz.
\item \textbf{$\alpha$ modelini OCUDU ölçümüyle karşılaştır.} Veri: \kod{ogretici/veri/ocudu\_prs\_sweep.csv} (PDSCH, logda \texttt{prb=[0,96)} olarak görülüyor, yani 96 PRB). Her PRS satırı için model kaybı:
$$\text{kayıp}_{model}=\frac{1}{T_{PRS}}\cdot\frac{\min(bw,96)}{96},\qquad \text{kayıp}_{ölçülen}=\frac{59{,}0-\text{dl\_brate}}{59{,}0}$$
Tablo: etiket, $T_{PRS}$, bw, model \%, ölçülen \%. Excel ya da Python.
\item \textbf{$\alpha$ modelini ns-3 ölçümüyle karşılaştır} (Aziz'in \kod{ns3\_bant.csv} dosyasından; Aziz üretince ya da \kod{ogretici/veri/ns3\_bant.csv}). $C(\alpha)/C_0$ ile $1-\alpha$ yan yana; $C_0$ = bant daraltmasız doygun throughput.
\item \textbf{Tip~1/2/3 yorumu.} \kod{docs/metrics.md} bölüm~4'teki OCUDU satırını kontrol et: ölçümde PRS 52\,PRB $\times$ 12 sembol; Tip~1 = PRS'in kapladığı kaynak. Aziz'in ve Fatih'in bulgularına göre (kartlarında) Tip~3 satırını güncelle.
\item \textbf{Rapor 2 taslağı:} Yöntem (üç araçta $\alpha$) ve Ön bulgular (3 tablo: adım 1, 2, 3). Ders raporu biçimi: birleşik rapor PDF'i (Taha'da).
\end{enumerate}

\begin{kontrol}
\textbf{Adım 2 (OCUDU):} p4 ($T=4$, 104\,PRB): model \%25,00, ölçülen \%25,76. p20: \%5,00 ve \%5,08. p10: \%10,00 ve \%10,17. p8: \%12,50 ve \%13,73. p4w52 ($T=4$, 52\,PRB): \%13,54 ve \%13,56. p80: \%1,25 ve \%1,69; p40: \%2,50 ve \%2,03.\\
\textbf{Adım 3 (ns-3, 1 band, 2 UE, doygun BE akışı):} $C_0=164{,}54$\,Mbps. $\alpha=0{,}20$: 133,41 ($C/C_0=0{,}811$); $\alpha=0{,}40$: 102,54 (0,623); $\alpha=0{,}50$: 77,06 (0,468). $\alpha=0{,}03$ (48,5\,MHz): 164,54 (oran 1,000); $\alpha=0{,}15$ (42,5\,MHz): 143,77 (0,874). $\alpha=0{,}80$ (10\,MHz): 0 paket alındı.\\
\textbf{Adım 1:} Case 2-4 (tek TRP, $\le\%10$, 1x1): $P_{md}$ 16,0\% ve 12,6--23,0\%; hedefler tutmuyor. Case 2-6 (tek TRP, $>\%10$, 1x1): 1 sonuçta $P_{md}$ 1,30\%, 3 sonuçta 12,30--37,80\%. Case 2-2 (çok TRP, $\le\%10$, 1x2/1x6/2x4): $P_{md}$ 1,27--4,67\%, tüm hedefler tutuyor.
\end{kontrol}

\begin{teslim}
(1) Baseline~2 sonuç tablosu; (2) OCUDU model--ölçüm tablosu ve 3 cümlelik yorum; (3) ns-3 model--ölçüm tablosu ve yorum; (4) güncellenmiş Tip~1/2/3 satırı; (5) Rapor 2 Yöntem ve Ön bulgular taslağı.
\end{teslim}

% ======================================================================= NIHAT
\clearpage
\section*{Nihat: Sionna (TR metrikleri ve varsayım karşılaştırması)}
\begin{neden}
TR 38.765 başarıyı $P_{md}\le\%5$, $P_{f1}\le\%5$, $P_{f2}\le\%5$ ve konum hatası $\le$10\,m (\%90) ile ölçüyor (§4.2). Elimizdeki radar bugüne kadar yalnızca $P_d$ ve RMSE verdi. Bu kart radarı standardın metrikleriyle okutuyor ve varsayımlarımızı standart ve literatürle yan yana koyuyor.
\end{neden}

\subsection*{Okuma}
\begingroup\sffamily\small
\begin{xltabular}{\textwidth}{@{}p{3.0cm}p{3.3cm}Xp{1.6cm}@{}}
\toprule
\textbf{Kaynak} & \textbf{Yer} & \textbf{Ne arayacaksın} & \textbf{Durum} \\
\midrule
Blandino vd. (arXiv 2505.24763), \kod{makaleler/} & III.A (denklem 9), III.C, III.D (denklem 15), IV.A (Tablo I), IV.B & III.A: statik clutter, PRS sembolleri boyunca ortalama çıkarılarak bastırılıyor. III.C: menzil için IDFT. III.D: tespit istatistiği = menzil tepesi / menzil kutularının ortalaması (PAR); eşik IV.B'de $P_{FA}$ ve $P_D$ eğrilerinden seçilmiş (3,4\,dB). IV.A: 30\,GHz, 120\,kHz, 66\,RB, comb $K=4$, 4 PRS sembolü, 256 PRS, 32$\times$32 dizi, EIRP 75\,dBm/100\,MHz, RCS $-12{,}81$\,dBsm, 4000 hedef atışı. & [Okundu] \\
Lin (arXiv 2608.11606) & IV.B & Monostatic sensing'in 5 adımı (sinyal yeniden oluşturma, kalibrasyon ve öz-girişim bastırma, gecikme-Doppler-açı işleme, clutter bastırma ve tespit, parametre kestirimi). & [Okundu] \\
Khosroshahi vd. (arXiv 2408.00667) & II.A & PRS: comb 2/4/6/12 ve sembol 1/2/4/6/12 seçenekleri; PRS boş RE bırakıyor ve bu menzil belirsizliği yaratıyor. & [Okundu] \\
Riihonen vd., \emph{Full-Duplex OFDM Radar With LTE and 5G NR Waveforms} (2019, IEEE TMTT; Consensus) & Özet & Verici--alıcı izolasyonu sınırlı olduğunda sorun esas olarak statik hedeflerin tespitinde; hareketli hedefler verici öz-girişimine karşı doğası gereği daha dayanıklı. & [Yalnızca özet] \\
\bottomrule
\end{xltabular}
\endgroup

\subsection*{Adımlar}
\begin{enumerate}[leftmargin=1.6em]
\item \textbf{Ön koşul:} \texttt{rt\_sehir.py} ve \texttt{radar\_takip.py} bir kez çalışmış olmalı (\kod{cikti/rt\_sonuclar.npz} ve \kod{cikti/kaynak\_basari.csv} oluşur). Python'a \texttt{pandas} kurulu olmalı.
\item \textbf{Lin'in 5 adımını kodla eşle.} \kod{radar\_takip.py} içinde her adımın hangi satırda olduğunu bul. Yazılacak tablo: \emph{Lin adımı | satır | var / yok}. Bizim modelde clutter yok (\texttt{radar()} yalnızca hedef ve gürültü üretiyor); bu satırı ``yok'' yaz ve III.A denklem~9'un ne yaptığını bir cümleyle yaz.
\item \textbf{TR metriklerini ölç:}
\begin{komut}
\begin{lstlisting}
$env:DRJIT_LIBLLVM_PATH = [Environment]::GetEnvironmentVariable('DRJIT_LIBLLVM_PATH','User')
cd ogretici\sionna_sehir
C:\Users\<kullanici>\venvs\ehb415\Scripts\python.exe tr_metrik_olc.py      # bizde 79 sn surdu
\end{lstlisting}
\end{komut}
Ubuntu'da: \texttt{\textasciitilde/venvs/ehb415/bin/python tr\_metrik\_olc.py}. Betik üç şey yazıyor: $P_{f1}$ (hedefsiz 3000 CPI), her menzil için $P_{md}\le\%5$ sağlayan en küçük $\alpha$, $X_{SI}=40$ ve 30\,dB için 300\,m'de $P_d$.
\item \textbf{Konum hatasının \%90 yüzdeliği.} \kod{radar\_takip.py}'deki \texttt{kaynak\_basari()} içinde \texttt{hr} (menzil hatası) ve \texttt{hv} (hız hatası) listeleri dolduruluyor. Her $(menzil, M)$ için \texttt{np.percentile(np.abs(hr), 90)} ve hız için aynısını yazdıran bir satır ekle. TR hedefleri: konum 10\,m, hız 5\,m/s (\%90).
\item \textbf{$X_{SI}$ sonucunu açıkla.} Betik çıktısındaki 3. bölüme bak. TR 38.765 Ek A (Tablo A-1, not~2): artık sızıntı ``$-94+X$\,dBm / 100\,MHz'' beyaz gürültü olarak modelleniyor, $X$ şirket seçimi. \kod{radar\_takip.py} satır~40'ta $X$ gürültüye ekleniyor. Neden $X$ değişince $P_d$ değişti? Cevabı bu iki yerden yaz.
\item \textbf{Varsayım karşılaştırma tablosu:} sütunlar \emph{bizim radar | Blandino IV.A | TR 38.765 Ek A}; satırlar taşıyıcı frekans, SCS, bant, verici gücü, hedef RCS, anten, tespit istatistiği, hedef sayısı. Bizim değerler \kod{radar\_takip.py} PARAM bölümünde (satır 23--32).
\end{enumerate}

\begin{kontrol}
\texttt{tr\_metrik\_olc.py} çıktısı (Windows, Python 3.11.9; tohum sabit):\\
$P_{f1}$: $M=14$ ($\alpha=\%10$): 2/3000 $\Rightarrow$ \%0,07; $M=70$ ($\alpha=\%50$): 5/3000 $\Rightarrow$ \%0,17.\\
$P_{md}\le\%5$ için en küçük $\alpha$: 100\,m: \%2,9 ($P_d=1{,}00$); 200\,m: \%20 (0,99); 300\,m: \%100 (0,975); 400\,m: taranan değerlerin hiçbirinde yok (en yüksek $P_d=0{,}38$).\\
300\,m, $\alpha=\%20/40/50/100$: $X_{SI}=40$\,dB: 0,18 / 0,56 / 0,73 / 0,96; $X_{SI}=30$\,dB: 1,00 / 1,00 / 1,00 / 1,00. (160 denemede standart hata en çok 0,04.)\\
Bu betiği Ubuntu'da (senin ortamında) çalıştırmadık; çıktı farklıysa farkı yaz.
\end{kontrol}

\begin{teslim}
(1) Lin--kod eşleme tablosu; (2) \texttt{tr\_metrik\_olc.py} çıktısı; (3) \%90 yüzdelik tablosu ve TR hedefiyle karşılaştırma; (4) $X_{SI}$ açıklaması (3--4 cümle); (5) varsayım karşılaştırma tablosu.
\end{teslim}

% ======================================================================= AZIZ
\clearpage
\section*{Aziz: ns-3 / 5G-LENA (yük, bant ve Sionna-RT durumu)}
\begin{neden}
ns-3 haberleşme bedelini ölçer. Dickerson modeli kapasitenin $(1-\alpha)$ ile azalacağını söylüyor; bu kart bunu ns-3'te ölçtürüyor ve sensing'i temsil etme yollarını (bant daraltma, ayrı trafik akışı) ölçümle karşılaştırtıyor. Raporun Yöntem bölümü bu karara bağlı.
\end{neden}

\subsection*{Okuma}
\begingroup\sffamily\small
\begin{xltabular}{\textwidth}{@{}p{3.0cm}p{3.3cm}Xp{1.6cm}@{}}
\toprule
\textbf{Kaynak} & \textbf{Yer} & \textbf{Ne arayacaksın} & \textbf{Durum} \\
\midrule
Ropitault vd. (arXiv 2508.04004), \kod{makaleler/} & III.A, III.B, III.C, IV.A & III.A: 5G-LENA'nın TR 38.901 kanal sınıfları. III.B: \texttt{TracesChannelModel} ve \texttt{TracesSpectrumPropagationLossModel}, dış ray tracer'ın çok yollu bileşenlerini (gecikme, genlik, faz, AoA/AoD) CSV'den okuyor. III.C: Sionna RT'den CSV üreten boru hattı (\texttt{wineslab/sionna-channel-generator}). IV.A: 28\,GHz, 100\,MHz, tek gNB--UE. & [Okundu] \\
Dickerson vd. (arXiv 2606.21677) & IV.D & Sabit politikalar: \%15 ayırınca throughput sınırı 0,85, \%3 ayırınca 0,97. ns-3'te iki politikayı bant daraltmayla taklit edeceksin. & [Okundu] \\
Capozzi vd., \emph{Downlink Packet Scheduling in LTE} (repo kökünde) & IV.A.2 (Round Robin), IV.B.1 (Maximum Throughput), IV.B.2 (Proportional Fair) & Her algoritmanın metriği; MT'nin zayıf kanallı kullanıcıya ne yaptığı. LTE makalesi; algoritma tanımları için. & [Okundu] \\
Koutlia vd. (arXiv 2205.03278) & 1 (giriş), 5.1 & 5G-LENA'nın 3GPP referans sonuçlarıyla kalibrasyonu, REM aracı. Rapor için: ``simülatör güvenilir mi'' cümlesinin kaynağı. & [Okundu] \\
Gautam vd., \emph{5G Configured Grant Scheduling for ISAC in Vehicular Networks} (2025, CNSM; Consensus) & Özet & ISAC sensing trafiği aperiyodik kabul edilip ns-3 5G-LENA CG modülünde Weibull gelişler ile modellenmiş. ``Sensing'i trafik olarak temsil'' yolunun literatürde kullanıldığını gösteriyor. & [Yalnızca özet] \\
\bottomrule
\end{xltabular}
\endgroup

\subsection*{Adımlar}
\begin{enumerate}[leftmargin=1.6em]
\item \textbf{Yük taraması} (Ubuntu, repo klasöründe):
\begin{komut}
\begin{lstlisting}
bash ogretici/betikler/ns3_yuk_sweep.sh      # cikti: /tmp/ns3_yuk.csv
\end{lstlisting}
\end{komut}
\texttt{cttc-nr-demo} iki akış kuruyor: akış~1 düşük gecikmeli (\texttt{packetSizeUll}, \texttt{lambdaUll}), akış~2 en iyi çaba (\texttt{lambdaBe}). Betik \texttt{--doubleOperationalBand=false} ile ikisini \textbf{aynı banda} koyar (varsayılan iki ayrı banttır) ve \texttt{lambdaBe}'yi 1000'den 40000'e çıkarır. Tablodan: BE akışının alınan throughput'u hangi gönderilen değerden sonra sabitleniyor?
\item \textbf{Bant taraması} (doygun yükte kapasiteyi ölçer):
\begin{komut}
\begin{lstlisting}
bash ogretici/betikler/ns3_bant_sweep.sh     # cikti: /tmp/ns3_bant.csv
\end{lstlisting}
\end{komut}
$\alpha=1-B/50$\,MHz. Tabloya sütun ekle: $C(\alpha)/C_0$ ve $1-\alpha$. İki ek ölçüm yap: \texttt{--bandwidthBand1=48500000} ($\alpha=\%3$) ve \texttt{42500000} ($\alpha=\%15$); Dickerson IV.D'deki iki politikaya karşılık gelir.
\item \textbf{İki soru (cevabı kendi ölçümünle ve kodla ver):} (a) 48,5\,MHz neden 50\,MHz ile aynı throughput'u veriyor? \kod{contrib/nr/examples/cttc-nr-demo.cc} içinde \texttt{bandwidthBand1}'in nerede kullanıldığını \texttt{grep} ile bul. (b) 10\,MHz'de neden hiç paket alınmıyor (\texttt{Rx Packets: 0})? Farklı değerlerle (örn.\ 15, 20\,MHz) sınırı bul.
\item \textbf{Sionna-RT örneğinin durumu.} 5G-LENA v5.1'de \texttt{cttc-nr-demo-sionna-rt} örneği var (\kod{contrib/nr/RELEASE\_NOTES.md} satır 275--278). \kod{contrib/nr/examples/CMakeLists.txt} içindeki \texttt{check\_deps(nr\_sionna\_deps ...)} bloğunu oku: örnek yalnızca \texttt{cppyy}, \texttt{sionna.rt}, \texttt{pybind11} Python paketleri bulunursa ve Windows değilse derleniyor. Bizim WSL'deki \texttt{python3} (3.12.3) için şunu çalıştır ve çıktıyı yaz:
\begin{komut}
\begin{lstlisting}
for m in cppyy sionna.rt pybind11; do python3 -c "import $m" 2>&1 | tail -1; done
cd ~/ehb415/ns-3-dev; ./ns3 run "cttc-nr-demo-sionna-rt --PrintHelp" 2>&1 | tail -3
\end{lstlisting}
\end{komut}
Paket kurmayı \textbf{deneme} (sistem Python'una dokunma); yalnızca durumu raporla.
\item \textbf{``Hangi yol'' notu (1 sayfa).} İki yolu bu kartın ölçümleriyle karşılaştır: (a) bant daraltma (adım~2), (b) sensing'i ayrı trafik akışı olarak ekleme (adım~1, \texttt{lambdaBe}). Her yol için: ns-3'te ne ölçüyor, hangi $\alpha$ tanımına karşılık geliyor, ölçümde ne görüldü. Önerini ve nedenini yaz. Perşembe toplantısında bu notu sen sunacaksın.
\end{enumerate}

\begin{kontrol}
\textbf{Adım 1} (\texttt{doubleOperationalBand=false}, 2 UE): akış~1 her \texttt{lambdaBe} için 10,23--10,24\,Mbps ve 0,28--0,32\,ms. Akış~2 (BE): gönderilen 10,24 $\rightarrow$ alınan 10,24; 51,20 $\rightarrow$ 51,18; 102,40 $\rightarrow$ 102,35; 204,80 $\rightarrow$ \textbf{164,54} (gecikme 59,35\,ms, kayıp \%19,66); 409,60 $\rightarrow$ \textbf{164,54} (104,47\,ms, kayıp \%59,83).\\
\textbf{Adım 2} (BE akışı alınan Mbps / gecikme ms): 50\,MHz 164,54 / 104,47; 48,5\,MHz 164,54 / 104,47; 42,5\,MHz 143,77 / 119,85; 40\,MHz 133,41 / 130,15; 30\,MHz 102,54 / 183,46; 25\,MHz 77,06 / 243,99; 10\,MHz 0 paket.\\
\textbf{Adım 4:} \texttt{cppyy}, \texttt{sionna.rt}, \texttt{pybind11}: üçü için de \texttt{ModuleNotFoundError}; \texttt{./ns3 run cttc-nr-demo-sionna-rt} bizde \texttt{Couldn't find the specified program} veriyor.\\
Varsayılan (\texttt{doubleOperationalBand=true}) iki ayrı banttır; bu durumda BE akışı 189,70\,Mbps'te doyuyor.
\end{kontrol}

\begin{teslim}
(1) Yük tablosu; (2) bant tablosu ($C/C_0$ ve $1-\alpha$ sütunlarıyla); (3) 3.\ adımdaki iki sorunun cevabı; (4) 4.\ adımın komut çıktıları; (5) 1 sayfa ``hangi yol'' notu; (6) 3 cümle: 5G-LENA'nın güvenilirliği (Koutlia).
\end{teslim}

% ======================================================================= FATIH
\clearpage
\section*{Fatih: OCUDU (PRS konumu, kod okuma, çok kullanıcı)}
\begin{neden}
OCUDU gerçek gNB kodu. Ölçtüğümüz bir bulgu var: PRS sembol sayısı DL throughput'unu değiştirmiyor. Bu kart nedenini kodda gösteriyor ve yeni bir gözlemi (PRS'in ofsetine göre UL'nin düşmesi) açık soru olarak sana bırakıyor. Cevaplar raporda ``OCUDU zamanlayıcısı PRS'e nasıl davranıyor'' bölümüne girer.
\end{neden}

\subsection*{Okuma}
\begingroup\sffamily\small
\begin{xltabular}{\textwidth}{@{}p{3.0cm}p{3.3cm}Xp{1.6cm}@{}}
\toprule
\textbf{Kaynak} & \textbf{Yer} & \textbf{Ne arayacaksın} & \textbf{Durum} \\
\midrule
Khosroshahi vd. (arXiv 2408.00667), \kod{makaleler/} & II.A, VI, VII & II.A: PRS comb 2/4/6/12, sembol 1/2/4/6/12; sensing ve iletişim sembolleri aynı RE'yi paylaşamaz. VI: PRS ile PDSCH arasında zaman ve frekans kaynak dağıtımı için optimizasyon problemi (denklem 27a--27e). VII: 28\,GHz, 120\,kHz için bulunan nokta: PRS'e 6 slot ve 12 PRB, kalan PDSCH'e. & [Okundu] \\
Lindenschmitt vd. (arXiv 2509.16917) & 3.3, 4.3, 4.4, 5 & 3.3 ve 4.3: radar işlemenin RU, DU ya da DU'ya bağlı RIC'te yapılması; gerçek zamanlılık, fronthaul yükü, maliyet ve güvenlik karşılaştırması. 4.4: dilimleme ile farklı SCS ve TDD düzenlerinin sensing kalitesine etkisi ve zamanlayıcının dengeyi nasıl kuracağı ``iyi anlaşılmış değil''. 5: ikinci bir ``sniffer RU'' ile standart değişikliği gerektirmeyen öneri. & [Okundu] \\
\midrule
OCUDU kodu (\kod{\textasciitilde/ehb415/ocudu/lib/scheduler/}) & \kod{common\_scheduling/prs\_scheduler.cpp} (satır 26--27, 193) & PRS PDU'sunun sembolleri \texttt{[symbol\_offset, symbol\_offset+nof\_symbols)}; satır~193'te bu semboller ve PRB'ler ızgarada \texttt{fill} ile dolu işaretleniyor. & [Okundu] \\
 & \kod{ue\_scheduling/intra\_slice\_scheduler.cpp} satır 865--912 & \texttt{update\_used\_dl\_vrbs}: dolu PRB'ler \texttt{common\_pdsch\_td\_resources()[0].symbols} aralığında taranıyor (satır 876--878); kod ``Implementation defined'' notu taşıyor. & [Okundu] \\
 & \kod{config/time\_domain\_resource\_helper.cpp} & Bu aralığın başlangıcı CORESET'in bittiği sembol, bitişi 14 (\texttt{generate\_dedicated\_pdsch\_td\_res\_list}, \texttt{calculate\_minimum\_pdsch\_symbol}). & [Okundu] \\
 & \kod{pdcch\_scheduling/pdcch\_slot\_resource\_allocator.cpp} satır 161--169 & PDCCH için ızgarada çakışma kontrolü (\texttt{collides}) ve \texttt{fill}. & [Okundu; satırlar arandı, tüm fonksiyon okunmadı] \\
\bottomrule
\end{xltabular}
\endgroup

\subsection*{Adımlar}
\begin{enumerate}[leftmargin=1.6em]
\item \textbf{Ofset deneyi} (Ubuntu, repo klasöründe; betik her koşuyu 50\,sn ile sınırlıyor). Betiğin yeni 5.\ argümanı PRS ilk sembol ofsetidir (varsayılan 2):
\begin{komut}
\begin{lstlisting}
bash ogretici/betikler/ocudu_prs_kos.sh e1 4 4  52 2      # 4 sembol,  52 PRB, ofset 2
bash ogretici/betikler/ocudu_prs_kos.sh e2 4 12 52 2      # 12 sembol, 52 PRB, ofset 2
bash ogretici/betikler/ocudu_prs_kos.sh e3 4 4  52 0      # ofset 0
bash ogretici/betikler/ocudu_prs_kos.sh e4 4 4  52 10     # ofset 10
bash ogretici/betikler/ocudu_prs_kos.sh e5 4 4  52 5      # ofset 5
\end{lstlisting}
\end{komut}
Her koşuda DL ve UL için \texttt{brate} ile \texttt{ok} sütunlarını oku (ilk satır ısınmadır; kararlı satırlara bak). e5 hata verirse \texttt{/tmp/out\_e5.txt}'deki hata satırını yaz.
\item \textbf{Soru 1 (kodla cevapla):} DL throughput ofset 2, 0 ve 10 için neden aynı? \kod{prs\_scheduler.cpp} satır~193 ve \kod{intra\_slice\_scheduler.cpp} satır 876--878'i okuyup bir paragraf yaz. İddialarını satır numaralarıyla destekle.
\item \textbf{Soru 2 (açık soru):} ofset 0'da DL değişmezken \textbf{UL neden düştü} (77,3 $\rightarrow$ 57,8\,Mbps, \texttt{ok} 1000 $\rightarrow$ 750)? Başlangıç noktası: \kod{pdcch\_slot\_resource\_allocator.cpp} satır 161--169. Cevabını \textbf{kanıtla}: log (\texttt{/tmp/out\_e3.txt} ya da \texttt{gnb} log düzeyini açarak) ya da kodda gördüğün satır. Kanıt bulamazsan ``bulamadım'' yaz ve nerelere baktığını sırala.
\item \textbf{Yapılandırma reddi:} e5'teki hata mesajını ve \kod{config/periodic\_resource\_sched\_validator.cpp}'yi oku. Hangi ofset değerlerinin kabul edildiğini bulmak için en az iki ofset daha dene (örn.\ 1, 7) ve tabloya \emph{kabul / ret} yaz. Ret nedenini mesajdan kopyala.
\item \textbf{Çok kullanıcılı deney:}
\begin{komut}
\begin{lstlisting}
bash ogretici/betikler/ocudu_cok_ue.sh u4   4 0 0      # 4 UE, PRS yok
bash ogretici/betikler/ocudu_cok_ue.sh u4p4 4 4 0      # 4 UE, PRS her 4 slotta
\end{lstlisting}
\end{komut}
\item \textbf{Mimari notu (yarım sayfa):} Lindenschmitt 3.3/4.3'e göre radar işlemenin RU, DU ya da RIC'te yapılmasının üç artı/eksisini yaz; Khosroshahi VII'deki nokta (PRS'e 6 slot ve 12 PRB) ile bizim ölçtüğümüz PRS kaybını (\texttt{ogretici/veri/ocudu\_prs\_sweep.csv}) yan yana koyan 3 cümle ekle.
\end{enumerate}

\begin{kontrol}
\textbf{Adım 1} (DL Mbps / ok; UL Mbps / ok): e1: 51 / 1000; 77,3 / 1000. e2: 51 / 1000; 77,3 / 1000. e3 (ofset 0): 51 / 1000; \textbf{57,8 / 750}. e4 (ofset 10): 51 / 1000; 77,3 / 1000. e5 (ofset 5): yapılandırma reddi: \texttt{Detected a collision between the NZP-CSI-RS resource (nzp-CSI-RS-ResourceId=4) and the DL-PRS resource (PRS Resource Set Id=0, PRS Resource Id=0) in slot 13 (period 20 slots)}. PRS yokken DL 59,0 ve UL 77,3.\\
\textbf{Adım 5} (toplam Mbps): 4 UE PRS yok: 60,0 (UE başına 15,0); 4 UE PRS her 4 slotta: 44,0 (UE başına 11,0). 1 UE: 59,0 ve 44,0; 2 UE: 60,0 ve 44,0; 8 UE: 60,8 ve 45,1.
\end{kontrol}

\begin{teslim}
(1) Ofset tablosu (DL, UL, kabul / ret); (2) Soru~1 paragrafı (satır numaralı); (3) Soru~2 cevabı ya da ``bulamadım'' ve baktığın yerler; (4) çok kullanıcılı tablo; (5) yarım sayfa mimari notu.
\end{teslim}

\end{document}
